% Written by Nghi (PL) — Project Lead
\section{Conclusion}
\label{sec:conclusion}

We presented a reproducibility-first empirical study evaluating the first-attempt,
zero-shot test-generation capability of \texttt{gpt-4o-mini-2024-07-18} on 50
standalone Java and Python functions with controlled cyclomatic complexity
(CC\,5--15).
Our protocol addresses a gap in the existing literature: no prior study had
simultaneously applied complexity filtering, zero-shot single-attempt evaluation,
and a paired search-based baseline within a pre-registered, fully logged design.

For RQ1, both language sub-populations reject the null hypothesis that median
branch coverage is at most 75\% (Java: $p = 0.000001$, 23/24 units above
threshold; Python: $p = 0.0004$, 19/22 units above threshold), demonstrating
that a single context-enriched prompt is sufficient to achieve high coverage on
CC-moderate functions without iterative repair.
For RQ2, GPT-generated tests substantially outperform Randoop 4.3.4 in mutation
score on all 24 paired Java units (median difference 62.1 percentage points,
Cliff's $\delta = 0.79$, $p = 6.6\times10^{-5}$), confirming that the semantic
understanding of the model provides a decisive advantage over random test
generation for functions with complex, domain-specific branching logic.

The primary failure mode identified is response truncation---three Python test
files ended with \texttt{SyntaxError: unterminated string literal}---suggesting
that increasing \texttt{max\_output\_tokens} to 4,096 would eliminate most
non-executable outputs without changing the generation strategy.
Python mutation score could not be evaluated due to a tool limitation in
\texttt{pytest-mutagen 1.3}, which we report as a construct-validity threat.

Our replication package, including the frozen manifest, prompt SHA-256 hashes,
raw API request/response logs, all metric outputs, and the pre-registered analysis
script, is publicly available at
\url{https://github.com/phamvohungnghi/SWT301-SE1944-Group6}.
Future work should extend the complexity range beyond CC\,15, incorporate an
automatic truncation-retry mechanism, and replicate the protocol with additional
LLMs to assess model-level variation under the same controlled design.
